\documentclass{article} % indique qu'on veut écrire un arcticle
\usepackage{xcolor}
\usepackage{graphicx}%pour importer des images
\usepackage{caption} % pour mattre une legende en dessous des images
\usepackage{subcaption}
\usepackage[utf8]{inputenc} % pour avoir les accent ect...
\usepackage[T1]{fontenc}
\usepackage[francais]{babel} % pour que ce soit écrit en francais
\usepackage{subfiles}% pour inclure d'autre fichier
\usepackage{hyperref} % Créer des liens et des signets
\hypersetup{
colorlinks=true, %colorise les liens
breaklinks=true, %permet le retour à la ligne dans les liens trop longs
urlcolor= blue, %couleur des hyperliens
linkcolor= black, %couleur des liens internes(dans la table des matiere)
citecolor=black, %couleur des références
}
\usepackage{enumitem}
\usepackage{float}
\usepackage{pifont}

\addtolength{\textwidth}{3cm} % Marges moins larges
\addtolength{\oddsidemargin}{-1.5cm} % Recentrage

\title{}
\author{} % indique l'auteur 
\date{}


\begin{document}
%-----------------------------------------------------------------------
\hrule
\begin{center}
\Huge {\textsc{CULTUR'ADVISOR - Développement d'une plateforme}}
\end{center}
\hrule
\vspace{10mm}
\begin{center}
\Large \textsc{Master 2 IMIS \\ Année 2022-2023}
\end{center}
\vspace{10mm}
\begin{figure}[H]
\hspace{12mm}
\includegraphics[scale = 1.2]{images/logo.png}
\end{figure}
\begin{figure}[H]
\hspace{30mm}
\includegraphics[scale = 0.3]{images/logo_dill.png}
\end{figure}
\begin{center}
\vfill{\Large \textsc{ Présenté par : Barodine Anaël, Mbelek-Nouga Paul-Dubien, TESSIER Adrien, Ververke Xavier \\ Encadré par : Tandjaoui Rosa et Hakim Aoudia}}
\end{center}
\newpage
%-----------------------------------------------------------------------
\vspace*{\stretch{1}}
\begin{center}
\Large\bfseries
\title -Remerciements
\end{center}
\hspace{4.5mm} 
Tout d'abord, nous tenons à remercier Rosa Tandjaoui, notre mentor, et Matthieu Exbrayat notre encadrant universitaire de nous avoir accompagnés et guidés pour ce projet.

Nous voulons aussi remercier toutes les personnes qui ont contribué au succès de ce projet et qui nous ont aidés lors de la rédaction de ce rapport.
\vspace*{\stretch{1}}
\newpage
%-----------------------------------------------------------------------
\tableofcontents
\newpage %aller a la page suivante
%-----------------------------------------------------------------------
\section{contexte du projet}
À la croisée de l’agenda culturel, du média, du magazine et du réseau social, CulturAdvisor est un espace de productions, de consultations et d’échanges autour de la culture. \\
Son objectif est de faciliter l’accès à la culture en aidant les acteurs à trouver de nouveaux publics, améliorer leur visibilité, augmenter leur fréquentation et les publics à trouver l’expérience culturelle qui leur ressemble, leur correspond.


\section{organisation de l'équipe}
Nous avons d'abord commencé par se répartir les différentes taches du projet.

\section{Projet}
Afin de réaliser ce projet, nous avons découper le projet en plusieurs taches.

\subsection{Mis en place du projet}
Premièrement, nous avons commencé par choisir un environnement de travail. \\
Nous avons choisit le langage Python et le framework Flask.

\subsection{Récupération des données}
Une fois le serveur Flask mis en place, il faut maintenant pouvoir récupérer les différentes données disponible.
\subsection{Normalisation des données}
Après avoir récupérer ces données, nous avons du les normaliser.

\subsection{requêtes}
Nous avons par la suite fait des requêtes...

\end{document}
